% !TEX TS-program = XeLaTeX
% !TEX encoding = UTF-8 Unicode

\chapter{讨论与适用性分析}
\label{chap:discussion}
\defaultfont

\section{方法结果的含义}
本文实验说明，在当前受试者和采集条件下，接收天线相位比与文件级统计特征能够形成具有类别区分能力的表示。该结果支持“人体姿态差异能够通过无线信道统计模式间接感知”这一研究思路，也说明在样本规模有限时，经过领域知识约束的统计表示可以与树集成模型形成有效组合。

然而，分类性能并不能单独证明模型学习到了唯一且稳定的姿态机理。无线信号同时受到人员、环境、设备和采集过程影响，模型可能利用多个相关因素完成分类。特征重要性能够反映模型内部使用哪些变量进行节点划分，但不等同于物理因果关系。要进一步解释信号变化与姿态之间的联系，需要通过控制变量实验、消融实验和跨条件重复测量排除其他可能因素。

\section{文件级统计表示的优势与局限}
文件级统计表示将长度可变的连续序列压缩为固定长度向量，能够降低单个异常帧和短时扰动对最终判断的影响。均值、分位数和标准差描述整体分布，差分统计描述时间变化强度，正弦和余弦统计保留周期变量的方向信息。对于持续时间较长、类别相对稳定的姿态状态，这种表示与研究对象具有较好的一致性。

统计聚合的主要局限是丢失精细时间顺序。两个具有不同变化过程的序列可能产生相近的分布统计量，短暂姿态转换也可能被完整文件的总体分布稀释。因此，该方法更适合离线文件级状态判断，不应在缺少验证的情况下直接用于快速动作分段或实时连续姿态追踪。若研究目标转向动态过程，需要结合时序编码、状态平滑或分段检测方法保留更多时间结构。

\section{人员与环境泛化}
无线信道具有明显的场景依赖性。受试者体型和动作习惯会改变人体反射特征，人与收发设备的距离和朝向会改变主要传播路径，家具移动和人员走动也会改变多径结构。因此，在相同人员和相近环境中取得的结果，不能直接代表陌生人员或新环境中的性能。

跨人员评价可以采用留一受试者验证或分组交叉验证，使测试人员完全不出现在训练阶段。跨环境评价需要在不同房间、收发位置和时间条件下采集独立数据；跨设备评价则需要考虑天线数量、子载波配置和射频链路差异。当输入维度或硬件特性改变时，原有特征定义可能无法直接迁移，需要重新建立链路对应关系或设计设备无关表示。

提高泛化能力可以从数据和方法两个方向开展。数据层面应增加人员、环境和采集时间的覆盖，并随机化姿态采集顺序；方法层面可以研究归一化、稳健特征选择、域适应和不确定性估计。无论采用何种方法，都应通过独立域测试验证，而不能仅根据训练集增强或模型复杂度推断泛化性能。

\section{潜在混杂因素}
姿态标签之外的因素若与类别同步变化，可能形成混杂。例如，不同姿态按照固定顺序采集时，设备温度、人员疲劳或环境变化可能与类别产生相关；不同类别若在不同位置或时间段采集，模型也可能利用位置和时段特征完成分类。混杂因素不会必然导致低准确率，反而可能产生较高但不可迁移的测试结果。

控制混杂需要在采集设计阶段随机化类别顺序、统一姿态持续时间并记录人员位置和环境状态。在分析阶段，可以按照人员、时间或采集批次进行分组评价，检查性能是否依赖某一批数据。对于错误样本，还应回看信号质量和采集记录，判断其是否与异常帧、姿态保持不稳定或环境扰动有关。

\section{辅助感知与医学应用边界}
本文模型的输入是无线信道状态信息，输出是预先定义的姿态类别。该过程没有直接测量脊柱曲率、椎体结构或其他临床指标，因此研究结果应解释为姿态辅助感知，而不是疾病筛查或诊断性能。即使分类标签使用了与脊柱形态相关的名称，模型仍然只是在当前标注体系下识别无线信号模式。

若未来面向健康监测场景，需要将技术有效性与临床有效性分开验证。技术有效性关注系统能否在不同人员和环境下稳定记录并识别姿态；临床有效性则需要专业医学标准、合适的受试者群体和规范的研究设计，验证系统输出是否与具有临床意义的指标相关。在完成这类验证之前，系统界面和论文表述都应避免使用替代诊断、确认疾病或保证健康结果等结论。

\section{隐私、伦理与数据治理}
\enlargethispage{\baselineskip}
无线感知不直接采集人体图像，能够降低面部和室内视觉信息暴露的风险，但 CSI 数据仍可能间接反映人员存在、活动规律和行为状态。因此，隐私保护不能只依据数据是否为图像判断。研究采集应取得知情同意，说明数据用途、保存期限和访问范围，并尽量减少与研究目标无关的身份信息。

模型部署还应考虑输出的可解释性和误判影响。对于健康相关提示，低置信度或输入质量不足时应允许系统拒绝判断，并向使用者说明结果的不确定性。数据访问、模型更新和结果记录需要建立权限与审计机制，避免无线感知数据被用于未经授权的人员监控。

\section{后续验证路线}
后续工作应按照由受控条件到开放条件的顺序推进。首先，可在相同人员和环境下开展不同时间的重复采集，评价结果随时间的稳定性；其次，采用按受试者分组的评价协议验证跨人员泛化；再次，在不同房间、收发位置和设备条件下建立独立测试集，分析环境和硬件变化的影响；最后，在完成离线稳定性验证后，再研究连续数据流中的在线分段、质量控制和结果平滑。

在方法分析方面，可以开展特征消融和参数敏感性实验。通过分别移除分布统计、差分统计和周期统计，可以判断各类特征的相对贡献；通过改变窗口长度、步长和树数量，可以观察模型对关键配置的敏感程度。所有比较应使用相同的数据划分和评价脚本，并报告重复实验的均值、波动或置信区间。只有实际完成这些实验后，才能对特征贡献和参数稳定性作出定量结论。

\section{本章小结}
本章讨论了文件级统计表示的适用条件、人员与环境泛化、潜在混杂因素、医学应用边界以及隐私伦理问题。当前结果能够支持方法在给定条件下的可行性，但不能代替跨人员、跨环境和临床有效性验证。后续研究需要通过分组评价、独立域测试、消融实验和不确定性处理逐步扩大证据范围。
